\documentclass{article}
\usepackage[utf8]{inputenc}
\usepackage[margin=1in]{geometry}
\usepackage{amsmath,amsthm,amsfonts,latexsym,amssymb,xcolor}
\usepackage{fancyhdr}
\usepackage{graphicx}
\usepackage{setspace}
\usepackage{hyperref}
\usepackage{csquotes}
\usepackage{booktabs} % For \toprule in table env
\usepackage{float}
\usepackage{tikz}
\usetikzlibrary{shapes, arrows, positioning, calc, graphs,intersections}
\usepackage{pgfplots}
% Define a custom style for the axes
\pgfplotsset{every axis/.append style={
    axis lines=middle, % No grid or box surrounding axes
    ticks=none, % Hide tick marks and labels
}}
\usepgfplotslibrary{fillbetween}
\pgfplotsset{compat=1.18}
\usepackage{imakeidx}
\makeindex
\usepackage[skip=\medskipamount]{parskip}
\pagestyle{fancy}
\setlength{\headheight}{25pt}
\setlength{\parindent}{0pt}
% Defining shade color for tikz graphic
\definecolor{shadecolor}{rgb}{0.8,0.8,0.8}

\newcommand{\interior}{\text{int}}
\newcommand{\cl}{\text{cl}}
\newcommand{\bd}{\text{bd}}
\newcommand{\R}{\mathbb{R}}
\newcommand{\N}{\mathbb{N}}
\newcommand{\Q}{\mathbb{Q}}
\newcommand{\Z}{\mathbb{Z}}
\newcommand{\toto}{\rightrightarrows}
\newcommand{\norm}[1]{\lVert #1 \rVert}
\newcommand{\sol}{\bigskip \bigskip \emph{Solution.} \bigskip}
\newcommand{\pf}{\bigskip \bigskip \emph{Proof.} \bigskip}
\newcommand{\keyword}[1]{\textit{\textbf{#1}}\index{#1}}

% argmax math operator
\DeclareMathOperator*{\argmax}{arg\,max}

% Using amsthm package
% Define my theorem/prop styling
\newtheoremstyle{mythm}
  {5pt}   % ABOVESPACE
  {5pt}   % BELOWSPACE
  {\itshape\onehalfspacing}  % BODYFONT
  {0pt}       % INDENT (empty value is the same as 0pt)
  {\bfseries} % HEADFONT
  {.}         % HEADPUNCT
  {5pt plus 1pt minus 1pt} % HEADSPACE
  {}       
\theoremstyle{mythm}
\newtheorem{theorem}{Theorem}
\newtheorem{corollary}{Corollary}
\newtheorem{lemma}{Lemma}
\newtheorem{claim}{Claim}
\newtheorem{proposition}{Proposition}
\newtheorem{assumption}{Assumption}

% Define my def/remark/ex styling
\newtheoremstyle{mydef}
  {5pt}   % ABOVESPACE
  {5pt}   % BELOWSPACE
  {\onehalfspacing}  % BODYFONT
  {0pt}       % INDENT (empty value is the same as 0pt)
  {\bfseries} % HEADFONT
  {.}         % HEADPUNCT
  {5pt plus 1pt minus 1pt} % HEADSPACE
  {}
\theoremstyle{mydef}
\newtheorem{definition}{Definition}
\newtheorem{remark}{Remark}
\newtheorem{ex}{Example}

\title{Lecture 4: Differentiation, IFT, Unconstrained Optimization}
\author{Originally written by Mauricio Caceres Bravo\\Revised by Cole Davis\\Webpage: \href{https://cj-davis99.github.io/math-camp.html}{cj-davis99.github.io}}
\date{August 16, 2024}

\onehalfspacing

\begin{document}
\lhead{Math Camp \\ August 16, 2024} 
\rhead{Lecture 4 \\ \leftmark}
\maketitle
\tableofcontents
\newpage
\section*{Notation}
\begin{itemize}
    \item $\forall$ translates to ``for all''
    \item $\exists$ translates to ``there exists''
    \item $\N = \{1,2, \hdots\}$ is the set of natural numbers
    \item $\Z = \{\hdots, -1, 0, 1, \hdots \}$ is the set of integers
    \item $\Q = \left\{ p/q : p \in \Z \text{ and } q \in \Z\backslash\{0\} \right\}$ is the set of rational numbers
    \item $\R$ is the set of real numbers
    \item If $S$ is a set and $n\in \N$, then $S^n$ is the $n^{\text{th}}$ order Cartesian product of $S$. E.g., $S^2 = S \times S$ 
    \item For any $\varepsilon>0$, $B_\varepsilon(x)$ is the Euclidean ball around $x$ with radius $\varepsilon$
    \item Unless otherwise specified, $d(x, y)$ is a metric on the contextual set $x,y$ belong to
    \item $\cl(S)$ is the closure of the set $S$
\end{itemize}
\newpage
\section{Differentiation}
\label{sub:differentiation}

A small warning to the reader: moving forward, we will state many results without proof. We do this for a few reasons:
\begin{itemize}
    \item[(a)] The results themselves are often much more useful than the proof of the results (e.g., we often just take the differentiability of a function for granted in the first-year courses).
    \item[(b)] The syllabus lists some textbooks that prove the results mentioned in this section if you are interested.
    \item[(c)] It is expected that the majority of incoming graduate students have quite a bit of exposure to this material already through calculus, linear algebra, differential equation, and analysis courses. 
\end{itemize}

\subsection{Single-Variable Calculus}

\begin{definition}
  Let $D\subseteq \R$, let $f: D \to \mathbb{R}$ be a function, and let $a\in D$ be a limit point\footnote{
      Recall that $x$ is a \keyword{limit point} of a set $D$ if $\forall \varepsilon > 0$, $B_\varepsilon(x) \cap D \neq \varnothing$.
  } of $D$. We say that $f$ is \keyword{differentiable} at $a$ if
  \[
    \lim_{x \to a} \dfrac{f(x) - f(a)}{x - a}
  \]
  converges to some real number $L$. We write $f'(a) = \frac{df}{dx}(a) = L$. If $f$ is differentiable $\forall a\in D$, then we say $f$ is differentiable on $D$.
\end{definition}

If $D$ is an interval in $\R$ (which it commonly is), then every point in $D$ is a limit point of $D$, and we don't need to worry about specifying that the point in which we want to evaluate a derivative is a limit point. 

Intuitively, differentiation yields the slope of a function at a point:

\begin{figure}[H]
  \centering
  \begin{tikzpicture}
    \begin{axis}[name=plot1
      %,title=
      ,width=15cm
      ,height=9cm
      ,ymin=1
      ,ymax=4.5
      ,xmin=0
      ,xmax=4.5
      ,domain=0:4
      ,ylabel=$f(x)$
      ,xlabel=$x$
    ]
    \addplot [smooth] {4 * x - x^2};
    \addplot [only marks, mark=*] coordinates {(2, 4)};
    \addplot [only marks, mark=*, mark size=1pt] coordinates {
      (1,    3)
      (1.5,  3.75)
      (1.75, 3.9375)
    };
    \addplot [only marks, mark=*, mark size=1pt] coordinates {
      (3,    3)
      (2.5,  3.75)
      (2.25, 3.9375)
    };
    \addplot [smooth, domain=0.55:3.45] {4}
      node[above] at (axis cs:3.15, 4) {$f^\prime(a)$};

    % lim x -> a-
    \draw [->, >=latex] (axis cs:1.1, 1.125) -- (axis cs:1.5, 1.125);
    \draw [->, >=latex] (axis cs:1.5, 1.125) -- (axis cs:1.75,1.125);
    \draw [->, >=latex] (axis cs:1.75,1.125) -- (axis cs:1.9, 1.125);
    \node[above] at (axis cs:1.0, 1) {$x$};
    \node[above] at (axis cs:2.0, 1) {$a$};
    \addplot [smooth, dashed, domain=0.50:2] {4 + 1.0  * (x - 2)};
    \addplot [smooth, dashed, domain=0.75:2] {4 + 0.5  * (x - 2)};
    \addplot [smooth, dashed, domain=1.00:2] {4 + 0.25 * (x - 2)};

    % lim x -> a+
    \draw [->, >=latex] (axis cs:2.9, 1.125) -- (axis cs:2.5, 1.125);
    \draw [->, >=latex] (axis cs:2.5, 1.125) -- (axis cs:2.25,1.125);
    \draw [->, >=latex] (axis cs:2.25,1.125) -- (axis cs:2.1, 1.125);
    \node[above] at (axis cs:3.0, 1) {$x$};
    \addplot [smooth, dashed, domain=2:3.50] {4 - 1.0  * (x - 2)};
    \addplot [smooth, dashed, domain=2:3.25] {4 - 0.5  * (x - 2)};
    \addplot [smooth, dashed, domain=2:3.00] {4 - 0.25 * (x - 2)};

    \end{axis}
  \end{tikzpicture}
  \caption{Graphical representation of a derivative}
  \label{fig:graphical_representation_of_a_derivative}
\end{figure}

\begin{theorem}
  If $f: D \to \mathbb{R}$ is differentiable at $a \in D$ then $f$ is also continuous at $a$.
\end{theorem}

\begin{proof}
  We want to show that $\forall \varepsilon > 0, ~ \exists \delta > 0$ s.t.
  \[
    |x - a| < \delta \implies |f(x) - f(a)| < \varepsilon
  \]

  Let $\varepsilon > 0$. Since $f$ is differentiable at $a$, we know we can find some $\widetilde{\delta} > 0$ s.t.
  \[
    0 < |x - a| < \widetilde{\delta} \implies \left|\dfrac{f(x) - f(a)}{x - a} - L \right| < \varepsilon
  \]
  


  where $L = f'(a)$. That is, we know the derivative exists and it is equal to $L$. Then we can write 

  \[
    \dfrac{|f(x) - f(a)|}{\widetilde{\delta}}
    < \dfrac{|f(x) - f(a)|}{|x - a|}
    = \left| \dfrac{f(x) - f(a)}{x - a} - L + L \right|
    \le \left| \dfrac{f(x) - f(a)}{x - a} - L \right| + |L|
    < \varepsilon + |L|
  \]
  
  Hence whenever $|x - a| < \widetilde{\delta}$ we get
  \[
    |f(x) - f(a)| < \left(\varepsilon + |L|\right) \widetilde{\delta}
  \]

  If we can find $\delta \le \widetilde{\delta}$ s.t. $\left(\varepsilon + |L|\right)\delta < \varepsilon$ then we'd be done. Take any $\delta$ such that $0 < \delta < \min\left(\tilde{\delta}, \dfrac{\varepsilon}{\varepsilon + |L|}\right)$. Then
  \[
    |x - a| < \delta < \widetilde{\delta}
    \implies
    |f(x) - f(a)|
    <
    \delta(\varepsilon + |L|)
    <
    \left(\varepsilon + |L|\right) \dfrac{\varepsilon}{\varepsilon + |L|}
    = \varepsilon
  \]
\end{proof}

\begin{theorem}\label{thm:diff_techniques}
  Let $D\subseteq \R$, let $a \in D$ be a limit point of $D$, and let $f: D \to \mathbb{R}$ and $g: D \to \mathbb{R}$ be differentiable at $a$.
  \begin{itemize}
    \item If $c\in\R$ is a constant, then $\dfrac{d [cf]}{dx} (a) = c f^\prime(a)$.
    \item $\dfrac{d[f + g]}{dx} (a) = f^\prime(a) + g^\prime(a)$.
    \item \keyword{Product rule}: $\dfrac{d[fg]}{dx} (a) = f^\prime(a) g(a) + f(a) g^\prime(a)$.
    \item \keyword{Power rule}: If $f(x) = x^k$ for some $k\in \R\setminus\{0\}$, then $\dfrac{df}{dx} (a) = k a^{k - 1}$.
    \item \keyword{Quotient rule}: If $g(a)\neq 0$, then $\dfrac{d}{dx} \dfrac{f(x)}{g(x)} = \dfrac{f^\prime(x) g(x) - f(x) g^\prime(x)}{g(x)^2}$.
  \end{itemize}
\end{theorem}

We state the final common derivative technique, the \keyword{chain rule}, as a separate theorem to signify the importance of the derivative technique and also the fact that its proof (which we do not include) is not as straightforward as the proofs in Theorem~\ref{thm:diff_techniques}:
\begin{theorem}
    Suppose $D, E \subseteq \R$, $f:D\to \R$ and $g:E\to\R$ are functions, $g(E)\subseteq D$, $g$ is differentiable at $a$, and $f$ is differentiable at $g(a)$. Then 
    \begin{equation*}
        \frac{d[f\circ g]}{dx}(a) = \frac{d[f(g)]}{dx}(a) = f'(g(a))g'(a)
    \end{equation*}
\end{theorem}

Some useful elementary function derivatives:
\begin{align*}
  \dfrac{d}{dx} e^x & = e^x  \\
  \dfrac{d}{dx} \ln(x) & = \dfrac{1}{x} \\
  \dfrac{d}{dx} \sin(x) & = \cos(x) \\
  \dfrac{d}{dx} \cos(x) & = -\sin(x)
\end{align*}

\begin{definition}
  A function $f: D \to \mathbb{R}$ is \keyword{continuously differentiable} if $f^\prime$ is continuous.
\end{definition}

\begin{ex}
  $f(x) = x^2$ is continuously differentiable since $f^\prime(x) = 2x$ is continuous. However,
  \begin{align*}
    f(x) = \begin{cases}
      x^2 \sin(1/x) & x \ne 0 \\
      0 & x = 0
    \end{cases}
  \end{align*}

  is \textit{not} continuously differentiable. In particular, for $x \ne 0$,
  \begin{align*}
    f^\prime(x)
    =
    2 x \cos(1/x) - x^2 \dfrac{1}{x^2} \cos(1/x)
    =
    2 x \cos(1/x) - \cos(1/x)
  \end{align*}

  and for $x = 0$ the derivative is $0$:
  \begin{align*}
    \lim_{x \to 0}
    \dfrac{f(x) - f(0)}{x - 0}
    =
    \lim_{x \to 0}
    \dfrac{x^2 \sin(1/x)}{x}
    =
    \lim_{x \to 0}
    x \sin(1/x)
  \end{align*}

  With $\sin(1/x)$ bounded, $x \to 0 \implies x \sin(x) \to 0$. However, the derivative itself is not continuous  at $0$. Note that while $2x \cos(1/x) \xrightarrow{x \to 0} 0$, $- \cos(1/x)$ does not have a limit, so the derivative does not have a limit as $x \to 0$ either, meaning it cannot be continuous.
\end{ex}

\begin{theorem}[L'H\^opital's rule]\label{thm:lecture4_lhopital}
Let $f$ and $g$ be continuous functions on $[a,b]$ and differentiable on $(a,b)$. Suppose that $c\in[a,b]$ and that $f(c) = g(c) = 0$. Suppose also that for some $\varepsilon > 0$, $g'(x) \neq 0$ for all $x\in \left([a,b] \cap B_\varepsilon(c)\right) \setminus \{c\}$. If 
\begin{equation*}
    \lim_{x\to c} \frac{f'(x)}{g'(x)} = L
\end{equation*}
then 
\begin{equation*}
    \lim_{x\to c} \frac{f(x)}{g(x)} = L
\end{equation*}
\end{theorem}

If $f, g$ are differentiable at $c$ (which is not actually required by the Theorem~\ref{thm:lecture4_lhopital}) then the intuition can be made plain, as this implies $f(c) = g(c) = 0$, and
\begin{align*}
    \lim_{x \to c} \dfrac{f(x)}{g(x)}
    =
    \lim_{x \to c} \dfrac{f(x) - 0}{g(x) - 0} \cdot \dfrac{x - c}{x - c}
    =
    \lim_{x \to c} \dfrac{(f(x) - f(c)) / (x - c)}{(g(x) - g(c)) / (x - c)}
    =
    \dfrac{f^{\prime}(c)}{g^{\prime}(c)}
\end{align*}

\subsubsection{Mean Value Theorem (MVT)}
\label{ssub:mean_value_theorem_mvt_}

We now move onto what is one of the most crucial results in single variable calculus: the Mean Value Theorem. 

\begin{theorem}[Mean Value Theorem]\label{thm:lecture4_mvt}
  Let $f: [a, b] \to \mathbb{R}$ be a continuous function on the closed interval $[a, b]$ and differentiable on the open interval $(a, b)$. Then $\exists c \in (a, b)$ s.t.
  \[
    f^\prime(c) = \dfrac{f(b) - f(a)}{b - a}
  \]
\end{theorem}

We will prove an equivalent, albeit perhaps conceptually easier, version of this:
\begin{theorem}[Rolle's Theorem]\label{thm:lecture4_rolles_theorem}
  Let $g: [a, b] \to \mathbb{R}$ be continuous on $[a, b]$ and differentiable on $(a, b)$ with $g(a) = g(b)$. Then $\exists c \in (a, b)$ s.t.
  \[
    g^\prime(c) = 0
  \]
\end{theorem}

\begin{figure}[!ht]
  \centering
  \begin{tikzpicture}
    \begin{axis}[name=plot1
      ,title={Mean Value Theorem}
      ,width=9cm
      ,height=8cm
      ,ymin=0
      ,ymax=3
      ,xmin=0
      ,xmax=2
      ,domain=0.25:1.75
      %,ylabel=$y$
      %,xlabel=$x$
    ]
    \addplot [smooth] {1 + (x - 0.5)^2}
      node[above] at (axis cs:0.25,   1.1625) {$f(a)$}
      node[above] at (axis cs:1.75,   2.6125) {$f(b)$}
      node[below] at (axis cs:1.0167, 1.167)  {$f(c)$};
    \addplot [smooth, dashed] {0.2167 + 1.033 * x}
      node[below] at (axis cs:1.75, 1.9) {$f^\prime(c)$};
    \addplot [only marks, mark=*] coordinates {
      (0.25,   1.0625)
      (1.75,   2.6125)
      (1.0167, 1.267)
    };
    \draw [dashed]
      (axis cs:0.25, 1.0625) --
      (axis cs:1.75, 2.6125);
    \end{axis}

    \begin{axis}[name=plot2
      ,title={Rolle's Theorem}
      ,at={($(plot1.east) + (0.5cm, 0)$)}
      ,anchor=west
      ,width=9cm
      ,height=8cm
      ,ymin=0
      ,ymax=4
      ,xmin=0
      ,xmax=4
      ,domain=0.25:3.75
      %,ylabel=$y$
      %,xlabel=$x$
    ]
    \addplot [smooth] {0.25 + (x - 2)^2}
      node[above] at (axis cs:0.25,   3.3125) {$f(a)$}
      node[above] at (axis cs:3.75,   3.3125) {$f(b)$}
      node[above] at (axis cs:2,      0.3)    {$f(c)$};
    \addplot [smooth, dashed] {0.25}
      node[below] at (axis cs:3.5, 0.6) {$f^\prime(c)$};
    \addplot [only marks, mark=*] coordinates {
      (0.25,   3.3125)
      (3.75,   3.3125)
      (2,      0.25)
    };
    \draw [dashed]
      (axis cs:0.25, 3.3125) --
      (axis cs:3.75, 3.3125);
    \end{axis}
  \end{tikzpicture}
  \caption{Graphical depiction of MVT}
  \label{fig:graphical_depiction_of_mvt}
\end{figure}

\begin{claim}
  \nameref{thm:lecture4_rolles_theorem} iff \nameref{thm:lecture4_mvt}.
\end{claim}

\begin{proof}
  The mean value theorem implies Rolle's theorem by definition. Simply note that by the MVT there is some $c$ s.t.
  \[
    g^\prime(c) = \dfrac{g(b) - g(a)}{b - a} = 0
  \]

  Now to show Rolle's theorem implies MVT, we only need a simple transformation:
  \[
    g(x) = \left(f(x) - f(a)\right) - \dfrac{f(b) - f(a)}{b - a} (x - a)
  \]

  (This is subtracting the line with slope $\dfrac{f(b) - f(a)}{b - a}$ that goes through $a$ from the function $f$, which is what we need to get it to ``flatten.'') Note that $g(a) = g(b) = 0$. Hence $\exists c$ s.t.
  \[
    g^\prime(c) = 0
  \]

  But we can see that
  \[
    g^\prime(c) = f^\prime(c) - \dfrac{f(b) - f(a)}{b - a} = 0
    \implies
    f^\prime(c) = \dfrac{f(b) - f(a)}{b - a}
  \]
\end{proof}

Now we show \nameref{thm:lecture4_rolles_theorem}.
\begin{proof}
    For Rolle's theorem we will use the extreme value theorem, and we know that $f$ is bounded and attains its $\sup$ and its $\inf$ in $[a, b]$. If the $\sup$ and the $\inf$ are \textit{both} at $\{a, b\}$, then because by assumption $f(a) = f(b)$, $f(x) = 0$ everywhere on $[a, b]$. Take any $c \in (a, b)$, and
    \[
        f^\prime(c)
        = \lim_{y \to c} \dfrac{f(y) - f(c)}{y - c}
        = \lim_{y \to c} \dfrac{0}{y - c}
        = \lim_{y \to c} 0 = 0
    \]

    Suppose then that either the $\sup$ or the $\inf$ occur at an interior point $c \in (a, b)$. Take the $\sup$ (WLOG; the $\inf$ is analogous or you can set $g = -f$). By assumption $f$ is differentiable, so $f^\prime(c)$ exists. That is, we know that
    \[
        f^\prime(c)
        = \lim_{y \to c} \dfrac{f(y) - f(c)}{y - c}
        = L
    \]

    Consider $y \to c^-$, that is $y$ approaching $c$ from the left. Since the $\sup$ is attained at $c$, $f(c) \ge f(y)$ for all $c > y$, which in turn gives
    \[
        \dfrac{f(y) - f(c)}{y - c} \ge 0
        \quad\quad
        \forall c > y
    \]

    Now take $y \to c^+$, that is $y$ approaching $c$ from the right. Again, since the $\sup$ is attained at $c$, $f(c) \ge f(y)$ for all $c < y$, which in turn gives
    \[
        \dfrac{f(y) - f(c)}{y - c} \le 0
        \quad\quad
        \forall c < y
    \]

    Which means that
    \[
      L_- = \lim_{y \to c^-} \dfrac{f(y) - f(c)}{y - c} \ge 0
      \quad\text{and}\quad
      L_+ = \lim_{y \to c^+} \dfrac{f(y) - f(c)}{y - c} \le 0
    \]

    Since the limit exists ($f$ is differentiable at $c$), we know $L = L_- = L_+$. Thus $0 \le L \le 0 \implies L = 0$.
    % Suppose not, then $L_- < 0$. But for every $\varepsilon > 0$ there is some $\delta$ s.t.
    % \[
    %     c^- - \delta < y < c^- < c^- + \delta
    %     \implies
    %     L_- - \varepsilon < \dfrac{f(y) - f(c)}{y - c} < L_- + \varepsilon
    % \]
    %
    % Take $\varepsilon$ s.t. $L_- < - \varepsilon < 0$. Then for some $\delta > 0$, we get
    % \[
    %     c^- - \delta < y < c^- < c^- + \delta
    %     \implies
    %     L_- - \varepsilon < \dfrac{f(y) - f(c)}{y - c} < L_- + \varepsilon < 0
    % \]
    %
    % But we have established that $y < c$ gives $\dfrac{f(y) - f(c)}{y - c} \ge 0$, a contradiction. By that same logic, we can write
    % \[
    %     L_+ = \lim_{y \to c^+} \dfrac{f(y) - f(c)}{y - c} \le 0
    % \]
    %
    % Since $f^\prime(c)$ exists, it must be the case that
    % \[
    %     f^\prime(c)
    %     = \lim_{y \to c^+} \dfrac{f(y) - f(c)}{y - c}
    %     = L_+
    %     = \lim_{y \to c^-} \dfrac{f(y) - f(c)}{y - c}
    %     = L_-
    %     = L
    % \]
    %
    % Now we can see that $L_+ \le 0$, $L_- \ge 0$, and $L_+ = L_- = L$  gives that $L = 0$. The proof for the case when the $\inf$ is interior is completely analogous.
\end{proof}

\begin{corollary}
  If $f$ is continuous on $[a, b]$ and differentiable on $(a, b)$ and obtains a local minimum or maximum at $c$, then $f^\prime(c) = 0$.
\end{corollary}

\begin{corollary}
  If $f$ is continuous on $[a, b]$ and differentiable on $(a, b)$ and $f^\prime(x) > 0$ for every $x \in (a, b)$ then $f$ is increasing on $(a, b)$. Conversely, if $f^\prime(x) < 0$ for every $x \in (a, b)$ then $f$ is decreasing on $(a, b)$.
\end{corollary}

\begin{theorem}[Cauchy's MVT]\label{thm:cauchy_mvt}
  Let $f, g: [a, b] \to \mathbb{R}$ be continuous functions on $[a, b]$ and differentiable on $(a, b)$. Then $\exists c \in (a, b)$ s.t.
  \[
    g^\prime(c) \left(f(b) - f(a)\right)
    =
    f^\prime(c) \left(g(b) - g(a)\right)
  \]
\end{theorem}

The MVT is a case where $g(x) = x$, the identity function.

\subsubsection{Taylor's Theorem}
Note that Theorem~\ref{thm:cauchy_mvt} is a generalization of the \nameref{thm:lecture4_mvt} (i.e., Theorem~\ref{thm:lecture4_mvt}); however, it's maybe not the most natural generalization of the MVT that one could formulate. Moreover, the scenarios where Theorem~\ref{thm:cauchy_mvt} is useful and Theorem~\ref{thm:lecture4_mvt} is not useful are scarce. 

A primary reason the MVT is so useful is that it provides us with an approximation of a differentiable function $f$. To see this, suppose $f:[a,b]\to\R$ meets all the requirements of Theorem~\ref{thm:lecture4_mvt} and we know the value of $f(a)$ and $f'(a)$, but we don't know the value of $f(b)$ (this is referred to as an interpolation problem in numerical analysis). By the MVT, we know there exists a $c\in (a,b)$ such that 
\begin{equation*}
    f(b) = f(a) + f'(c)(b-a)
\end{equation*}
We don't know $c$ (or $f'(c)$ for that matter), but if $b$ is sufficiently close to $a$, then it stands to reason that 
\begin{equation*}
    f(b) \approx f(a) + f'(a)(b-a)
\end{equation*}
More generally, it allows us to write 
\begin{equation*}
    f(x) \approx f(a) + f'(a)(x-a)
\end{equation*}
whenever $x\in [a,b]$.

A few natural questions are raised through this thought process:
\begin{itemize}
    \item How accurate is this approximation for a general $x\in (a,b)$? 
    \item Can we generate a more accurate approximation for a general $x\in (a,b)$? 
\end{itemize}
The following theorem provides a useful generalization of the MVT that addresses the questions above: 

\begin{theorem}[Taylor's theorem]
  Suppose $f:[a,b]\to \R$ is continuously differentiable $n+1$ times on $(a,b)$ and $x_0\in (a,b)$. Then for each $x\in[a,b]$, there exists a $c$ between $x$ and $x_0$ such that 
  \begin{equation*}
      f(x) = \underbrace{\sum_{k=0}^n \frac{f^{(k)}(x_0)}{k!} (x-x_0)^k}_{T_n(x)} + \underbrace{\frac{f^{(n+1)}(c)}{(n+1)!}(x-x_0)^{n+1}}_{R_n(x)}
  \end{equation*}
  where $T_n(x)$ is called the \keyword{$n^{\text{th}}$ order Taylor polynomial} and $R_n(x)$ is called the remainder or error. 
\end{theorem}

\begin{ex}
  \begin{itemize}
    \item The Taylor expansion of any polynomial is the polynomial itself. Consider $f(x) = x^2 + 3x$
      \begin{align*}
        f(x)
        & =
        \dfrac{f(1)}{0!} (x - 1)^0
        +
        \dfrac{f^\prime(1)}{1!} (x - 1)^1
        +
        \dfrac{f^{\prime\prime}(c)}{2!} (x - 1)^2
        \\
        & =
        4
        +
        5 (x - 1)
        +
        \dfrac{2}{2} (x - 1)^2
        \\
        & =
        4 + 5x - 5 + x^2 - 2x + 1
        =
        x^2 + 3x
      \end{align*}

      However, for lower-order Taylor expansions the theorem still applies. Visually:
      \begin{figure}[H]
        \centering
        \caption{Visualizing Taylor's Theorem}
        \label{fig:visualizing_taylor_s_theorem}
        \begin{tikzpicture}
          \begin{axis}[name=plot1
            ,title={Given $x = 5$ Recover Exact Value}
            ,width=9cm
            ,height=8cm
            ,ymin=0
            ,ymax=6
            ,xmin=0
            ,xmax=1.5
            ,domain=0:1.5
            ,ylabel={$f(x) = x^2 + 3x$}
            ,xlabel=$x$
            ,clip=false
          ]
          \addplot [smooth,domain=0:1.25] {x^2 + 3 * x};
          \addplot [smooth,domain=0.1:1.15] {4 + (2 * (3/4) + 3) * (x - 1)};
          \draw [dashed] (axis cs:0.5, 1.75) -- (axis cs:0.5, 0);
            \node[below] at (axis cs:0.5, 0) {$x = 0.5$};
          \draw [->, >=latex] (axis cs:0.5, 3.5) -- (axis cs:0.5, 1.95);
            \node[above] at (axis cs:0.5, 3.5) {$c = 0.75 \in (0.5, 1)$};
          \addplot [only marks, mark=*] coordinates {(0.5, 1.75)};
          \end{axis}

          \begin{axis}[name=plot2
            ,title={Approximate the function around $x = 1$}
            ,at={($(plot1.east) + (0.5cm, 0)$)}
            ,anchor=west
            ,width=9cm
            ,height=8cm
            ,ymin=0
            ,ymax=7
            ,xmin=0
            ,xmax=1.75
            ,domain=0:1.75
            ,ylabel={$f(x) = x^2 + 3x$}
            ,xlabel=$x$
            ,clip=false
          ]
          \addplot [smooth,domain=0:1.5] {x^2 + 3 * x};
          \addplot [smooth,domain=0.15:1.5] {4 + 5 * (x - 1)};
          \draw [dashed] (axis cs:1, 4) -- (axis cs:1, 0);
            \node[below] at (axis cs:1, 0) {$x = 1$};
          \addplot [only marks, mark=*] coordinates {(1, 4)};
          \draw [->, >=latex] (axis cs:1.35, 4.75) -- (axis cs:1.35, 5.65);
            \node[below] at (axis cs:1.35, 4.75) {$5x - 1$};
          \end{axis}
        \end{tikzpicture}
      \end{figure}

      The left and right figures correspond to:
      \begin{align*}
        f(x)
        & = \dfrac{f(1)}{0!} (x - 1)^0 + \dfrac{f^\prime(c)}{1!} (x - 1)^1
        \\
        & = 4 + (2c + 3) (x - 1)
        \\
        f(x)
        & \approx \dfrac{f(1)}{0!} (x - 1)^0 + \dfrac{f^\prime(1)}{1!} (x - 1)^1
        \\
        & = 4 + 5 (x - 1)
      \end{align*}

      We can see on the left that there is indeed some number $c$ s.t. Taylor's theorem holds. For our sample value of $x = 0.5$ we find $c = 0.75 \in (0.5, 1)$. On the right figure, on the other hand, we plot the \textit{approximation}. In this case, the function and the approximation are exactly equal at $1$, since the approximation at that point simplifies to $f(1)$. We can also see that around $x = 1$ the approximation is fairly good! However, farther away the error increases, as we would expect.

    \item One common Taylor approximation is for the logarithm. In particular, the first order Taylor expansion around  $1$:
      \begin{align*}
        \log(x) \approx \log(1) + \dfrac{1}{1} (x - 1) = x - 1
      \end{align*}

      Another version of this approximation is for $\log(1 + x)$ around $0$:
      \begin{align*}
        \log(1 + x) \approx \log(1) + \dfrac{1}{1 + 0} (x - 0) = x
      \end{align*}

      Take the relation:
      \begin{align*}
        \log Y = \alpha \log K + (1 - \alpha) \log L
      \end{align*}

      Using the approximation above, we can say that if $K$ increases by $10\%$, then $Y$ increases by $\alpha \log (1.1) \approx \alpha \cdot 0.1$, that is, $10\alpha \%$.
  \end{itemize}

  This second approximation is often used when dealing with percentage changes. You will hear the term \keyword{log-linearization}  thrown around, and this is what that's in reference to: Logarithms can be approximated as a percentage for small values.
\end{ex}

\subsection{Partial Derivatives}
\label{ssub:partial_derivatives}

\begin{definition}
  Let $f: S \subseteq \mathbb{R}^N \to \mathbb{R}^M$, $e = (e_1, \ldots, e_N)$ a standard basis of $\mathbb{R}^N$, $u = (u_1, \ldots, u_M)$ a standard basis for $\mathbb{R}^M$. $\forall x \in S, f(x) \in \mathbb{R}^M$, $f(x)$ is a linear combination of $u$ and some set of functions $\{f_1, \ldots, f_M\}$ s.t. $f_i: S \to \mathbb{R}$ with
  \[
    f(x) = \sum^{M}_{i = 1} f_i(x) u_i
  \]
\end{definition}

\begin{definition}
  The \keyword{partial derivative} of a function $f: S \subseteq \mathbb{R}^N \to \mathbb{R}^M$ at $x \in S$ is
  \[
    \dfrac{\partial}{\partial x_j} f_i(x) = \lim_{t \to 0} \dfrac{f_i(x + t \cdot e_j) - f_i(x)}{t}
  \]

  for $i = 1, \ldots, M$ and $j = 1, \ldots, N$.
\end{definition}

Note that partial derivatives needn't imply anything about the behavior of the function overall. Take, for instance, $f: \mathbb{R}^2 \to \mathbb{R}$ with
\[
  f(x, y) =
  \begin{cases}
    \dfrac{x^2y^4}{x^4 + y^8} & (x, y) \ne (0, 0) \\
    0 & (x, y) = (0, 0)
  \end{cases}
\]

The partial derivatives at $0$ are all $0$ (crucially, we take the limits one dimension at a time, so it's not that $(x, y) \to (0, 0)$, but rather $x \to 0$ and $y \to 0$ separately. Now take any $y_m \to 0$ and $x_m = y_m^2$ s.t. $y_m \ne 0$ for all $n$.
\[
  \lim_{n \to \infty} f(x_m, y_m) = \dfrac{1}{2} \ne 0 = f(0, 0)
\]

which means the function is not even continuous at $0$.

\begin{theorem}[Schwarz's\footnote{I have more commonly seen this referred to as Clairaut’s Theorem; however, according to \href{https://en.wikipedia.org/wiki/Symmetry_of_second_derivatives}{Wikipedia}, Schwarz published the first officially accepted proof of the result.} Theorem]
  Let $f: S \subseteq \mathbb{R}^N \to \mathbb{R}$; then
  \[
    \dfrac{\partial}{\partial x_i} \dfrac{\partial}{\partial x_j} f(x)
    =
    \dfrac{\partial}{\partial x_j} \dfrac{\partial}{\partial x_i} f(x)
  \]

  if all the partial derivatives exist and are continuous; that is, mixed partials are symmetric.
\end{theorem}

\begin{definition}
  The \keyword{gradient} of $f: S \subseteq \mathbb{R}^N \to \mathbb{R}$ at $x \in S$ is
  \[
    (\nabla f)(x) = \left[\begin{matrix}
      \dfrac{\partial}{\partial x_1} f(x) \\
      \vdots \\
      \dfrac{\partial}{\partial x_N} f(x)
    \end{matrix}\right]
  \]

  The gradient can also be denoted as $(Df)(x)$ or $D_x f(x)$.
\end{definition}

\begin{definition}
  The \keyword{Hessian} of $f: S \subseteq \mathbb{R}^N \to \mathbb{R}$ at $x \in S$ is
  \[
    (D^2 f)(x) = \left[\begin{matrix}
      \dfrac{\partial^2}{\partial x_1^2} f(x) & \cdots & \dfrac{\partial^2}{\partial x_N \partial x_1} f(x) \\
      \vdots & \ddots & \vdots \\
      \dfrac{\partial^2}{\partial x_1 \partial x_N} f(x) & \cdots & \dfrac{\partial^2}{\partial x_N^2} f(x) \\
    \end{matrix}\right]
  \]

  Where $D^2$ denotes the application of the $D$ operator twice. The Hessian can also be denoted $D_{x^\prime} D_x f(x)$ (this is analogous to the $d/dx^2$ operator for univatiate functions, where you can think $x^2$ for a vector is $x^\prime x$).
\end{definition}

Note: I am using $x^\prime$ to denote the transpose; alternatively this can be denoted $x^T$.
\begin{definition}
  Let $f: S \subseteq \mathbb{R}^N \to \mathbb{R}^M$ and $x \in S$; the \keyword{Jacobian} of $f$ is the $M \times N$ matrix with $i, j$ entry equal to $\dfrac{\partial}{\partial x_j} f_i(x)$, denoted $D_{x^\prime} f(x)$.
\end{definition}

\section{Implicit Function Theorem (IFT)}
\label{sub:implicit_function_theorem_ift_}

\subsection{Motivation}
\label{ssub:motivation}

Consider a function $f(x, y) = 0$. How do we characterize a function relating $x$ to $y$? We can write that for $f: \mathbb{R}^2 \to \mathbb{R}$ continuously differentiable on an open set $O$ with $f(x, y) = 0$, there exists some function $h$ s.t.
\[
  f(x, h(x)) = 0
  \quad \text{ and } \quad 
  \dfrac{dy}{dx} = -\dfrac{\partial f / \partial x}{\partial f / \partial y}
\]

One classic example is how to characterize the slope of a tangent line at a point $(x, y)$ of some circle of radius $r$ centered at $(0, 0)$. That is, $x^2 + y^2 = r^2$. We can write
\[
  f(x, y) = x^2 + y^2 - r^2 = 0
\]

So we know that for some $h(x)$, $f(x, h(x)) = 0$. Furthermore,
\[
  \dfrac{\partial f}{\partial x}
  = 2x
  \quad
  \dfrac{\partial f}{\partial y}
  = 2y
  \implies
  \dfrac{dy}{dx}
  = - \dfrac{x}{y}
\]

We will look at the general version of the theorem. We often work with a parameter space and a variable space, and we  want to express the variables in terms of the parameters (or with exogenous and endogenous variables, and we want to express the endogenous variables in terms of the exogenous variables). Take
\[
  (\theta_1, \ldots, \theta_N) = \theta \in \mathbb{R}^N
\]

to be the parameters (or exogenous) and
\[
  (x_1, \ldots, x_M) = x \in \mathbb{R}^M
\]

to be the variables (or endogenous). It's typically not common to have an explicit expression for the latter in terms of the former, but often we will encounter an implicit relation of the form
\[
  f(\theta, x) = 0
\]

For example, some system of equations
\begin{equation}
  \begin{array}{c}
    f_1(\theta, x) = 0 \\
    \vdots \\
    f_M(\theta, x) = 0
  \end{array}
  \nonumber
\end{equation}

The IFT gives a result we can apply to these types of problems.

\subsection{The IFT}
\label{ssub:the_ift}

\begin{theorem}[Implicit Function Theorem]

  Suppose $U\subseteq \R^{N+M}$ is open and $f:U\to \R^M$ is continuously differentiable on $U$. Fix a point $(\theta_0, x_0) \in U$ with $f(\theta_0, x_0) = 0 \in \R^M$. If $D_{x'}f(\theta_0, x_0)$ is non-singular (i.e. full-rank, or has a non-zero determinant) then there exists an open set $V\subseteq \R^N$ and a function $g:V\to \R^M$ such that 
  \begin{enumerate}
      \item $\theta_0 \in V$
      \item $g(\theta_0) = x_0$ 
      \item $f(\theta, g(\theta)) = 0$ for all $\theta \in V$
  \end{enumerate}
  Moreover, $g$ is continuously differentiable on $U$ and differentiating with respect to $\theta'$ yields
  \begin{align*}
      D_{\theta'}f(\theta, g(\theta)) + D_{x'} f(\theta, g(\theta)) D_{\theta'}g(\theta) = 0 \\
      D_{\theta'}g(\theta) = -\left[D_{x'} f(\theta, g(\theta))\right]^{-1}D_{\theta'}f(\theta, g(\theta))
  \end{align*}
\end{theorem}

\subsubsection{Example}
\label{ssub:example}

Take a simplified version of the IS-LM model
\begin{align*}
  Y   & = C + I + G   \\
  C   & = C(Y - T)    \\
  I   & = I(r)        \\
  M^S & = M^D(Y, r)
\end{align*}

with
\[
  0 < C^\prime(x) < 1
  \quad\quad
  I^\prime(r) < 0
  \quad\quad
  \dfrac{\partial M^D}{\partial Y} > 0
  \quad\quad
  \dfrac{\partial M^D}{\partial r} < 0
\]

National income must equal consumption plus investment (savings) plus government spending; consumption is some function of income minus taxes, the level of investment is determined by the interest rate, and money supply must equal money demand.  We have that
\begin{align*}
  Y - C(Y - T) - I(r) - G & = 0 \\
  M^S - M^D (Y, r)        & = 0
\end{align*}

which is the exact type of problem the IFT can help us solve. We have endogenous variables $x = (Y, r)$, national income and the interest rate, and exogenous variables $\theta = (M^S, G, T)$, money supply, government spending, and taxes. Hence
\begin{equation}
 \label{eq:is_lm_example}
  f(\theta, x)
  =
  \left[\begin{matrix}
    f_1(\theta, x) \\
    f_2(\theta, x)
  \end{matrix}\right]
  =
  \left[\begin{matrix}
    Y - C(Y - T) - I(r) - G \\
    M^S - M^D (Y, r)
  \end{matrix}\right]
  = 0
\end{equation}

Hence for some $g$, we can write
\begin{equation}
  \begin{array}{c}
    g(\theta)
    = \left[\begin{matrix}
      Y(M^S, G, T) \\
      r(M^S, G, T)
    \end{matrix}\right] \\[6pt]
    D_{\theta^\prime} f(\theta, g(\theta)) + D_{x^\prime} f(\theta, g(\theta)) D_{\theta^\prime} g(\theta) = 0
  \end{array}
  \nonumber
\end{equation}

(\textbf{Note}: \textbf{\textit{Stop here in class; the rest is a homework problem.}}) We have that
\begin{align*}
  D_{\theta^\prime} f(\theta, g(\theta))
  & =
  \left[\begin{matrix}
      0
    & -1
    & C^\prime(Y(\cdot) - T) \\[6pt]
      1
    & 0
    & 0
  \end{matrix}\right] \\[6pt]
  D_{x^\prime} f(\theta, g(\theta))
  & =
  \left[\begin{matrix}
      \dfrac{\partial f_1}{\partial Y} = 1 - C^\prime(Y(\cdot) - T)
    & \dfrac{\partial f_1}{\partial r} = -I^\prime(r(\cdot)) \\[6pt]
      \dfrac{\partial f_2}{\partial Y} = -\dfrac{\partial M^D}{\partial Y}
    & \dfrac{\partial f_2}{\partial r} = -\dfrac{\partial M^D}{\partial r}
  \end{matrix}\right] \\[6pt]
  \left[D_{x^\prime} f(\theta, g(\theta))\right]^{-1}
  & =
  \dfrac{1}{\text{det}(D_{x'}f(\theta,g(\theta)))}
  \left[\begin{matrix}
      - \dfrac{\partial M^D}{\partial r}
    & I^\prime(r(\cdot)) \\[6pt]
      \dfrac{\partial M^D}{\partial Y}
    & 1 - C^\prime(Y(\cdot) - T)
  \end{matrix}\right] \\[6pt]
  D\equiv \text{det}(D_{x'}f(\theta,g(\theta))) & =
  \underbrace{
   - \overbrace{\dfrac{\partial M^D}{\partial r}}^{< 0}
     \underbrace{\left(1 - C^\prime(Y(\cdot) - T\right)}_{> 0}
   }_{> 0}
  \underbrace{
   - \overbrace{I^\prime(r(\cdot))}^{< 0}
     \underbrace{\dfrac{\partial M^D}{\partial Y}}_{> 0}
   }_{> 0}
   \implies D > 0
\end{align*}

A non-zero determinant implies that the inverse exists. Hence we find that
\begin{align*}
  D_{\theta^\prime} g(\theta)
  & =
  \left[\begin{matrix}
      \dfrac{\partial Y}{\partial M^S}
    & \dfrac{\partial Y}{\partial G}
    & \dfrac{\partial Y}{\partial T} \\[6pt]
      \dfrac{\partial r}{\partial M^S}
    & \dfrac{\partial r}{\partial G}
    & \dfrac{\partial r}{\partial T}
  \end{matrix}\right]
  =
  - \dfrac{1}{D}
  \left[\begin{matrix}
        I^\prime(r(\cdot))
    &   \dfrac{\partial M^D}{\partial r}
    & - \dfrac{\partial M^D}{\partial r} C^\prime(Y(\cdot) - T)
    \\[6pt]
        1 - C^\prime(Y(\cdot) - T)
    & - \dfrac{\partial M^D}{\partial Y}
    &   \dfrac{\partial M^D}{\partial Y} C^\prime(Y(\cdot) - T)
    \\[6pt]
  \end{matrix}\right] \\
  & =
  - \dfrac{1}{D}
  \left[\begin{matrix}
      < 0
    & < 0
    & > 0 \\[6pt]
      > 0
    & < 0
    & > 0
  \end{matrix}\right]
  =
  \dfrac{1}{D}
  \left[\begin{matrix}
      > 0
    & > 0
    & < 0 \\[6pt]
      < 0
    & > 0
    & < 0
  \end{matrix}\right]
\end{align*}

Which means that for some $x = (Y, r), \theta = (M^S, G, T)$ that satisfies \eqref{eq:is_lm_example} there is some local neighborhood around $\theta$ where we can characterize the behavior of $(Y, r)$ with respect to each of the variables in $\theta$. In particular, income reacts positively to increases money supply or government spending but negatively to taxes, while the interest rate goes down with increases in the money supply or taxes but goes up with increases in government spending.

\section{Unconstrained Optimization}
\label{sub:unconstrained_optimization}


\newpage
\printindex
\end{document}